\documentclass[11pt]{article}
\usepackage{graphicx} % Required for inserting images
\usepackage{amsmath}
\usepackage{amssymb}
\usepackage{bm}
\usepackage{xcolor}
\usepackage{float}
\usepackage{listings}
\usepackage[mathscr]{eucal}
\usepackage[most]{tcolorbox}

\DeclareTColorBox{notebox}{ O{Note} O{blue} }{
  colback=#2!5!white,
  colframe=#2!75!black,
  fonttitle=\bfseries,
  title={#1}
}

\title{Radiation-Hydrodynamics Simulation Verification Suite}
\author{Tal Ramot}
\date{Exams week 3, 2026}

\begin{document}
\maketitle

\section{Setups for IMC and Rad-Hydro Simulations}

This section contains exact physical benchmarks to verify supersonic radiative transfer implementations. The parameters are calculated to isolate local thermodynamic equilibrium (LTE) supersonic diffusion behavior against known analytical profiles.

\subsection{Test 1 - Supersonic Henyey Wave (Power-Law Surface Drive)}
For the first verification case, we deploy a supersonic Henyey wave matching a gold-like material model at a high-energy characteristic scale of $1\,\text{HeV} = 0.1\,\text{KeV} \approx 1.160452 \times 10^6\,\text{K}$. At this scale, the material energy density dominates, and independent radiation energy density ($a_R T^4$) can be neglected in the material energy balance during non-LTE tracking.

\subsubsection{Material Model and Equation of State}
The Rosseland mean opacity and specific internal energy are defined using power-law dependencies on temperature $T$ under a constant spatially homogeneous material density $\rho_0$. 
We set the material exponents to $\alpha = 1.5$ and $\gamma = 1.6$, with a density coupling parameter $\mu = 0$ and $\lambda = 0$.

The macroscopic opacity matches the parameterization:
\begin{equation}
    k(T) = \rho_0 \kappa_R(T) = k_0 T^{-\alpha} = 10^4 \left(\frac{T}{\text{HeV}}\right)^{-1.5} \,\text{cm}^{-1}
\end{equation}
Thus, the opacity prefactor is $k_0 = 10^4\,\text{HeV}^{1.5}/\text{cm}$, which implies the face-averaging prefactor is:
\begin{equation}
    g = \frac{1}{\rho_0 k_0} = \frac{1}{10^4 \rho_0} \,\text{cm}\cdot\text{HeV}^{-1.5}
\end{equation}

The specific internal energy equation of state follows:
\begin{equation}
    e(T) = f T^{\gamma} = f T^{1.6}
\end{equation}
By matching the total energy density to the standard benchmark equation $u(T) = \rho_0 e(T) = u_0 T^b$ (where $b = \gamma = 1.6$), we select $u_0 = 10^{13}\,\text{erg}/(\text{cm}^3\cdot\text{HeV}^{1.6})$. The material prefactor input is therefore:
\begin{equation}
    f = \frac{u_0}{\rho_0} = \frac{10^{13}}{\rho_0} \,\text{erg}/(\text{g}\cdot\text{HeV}^{1.6})
\end{equation}

\subsubsection{Boundary Conditions and Driving Parameters}
The problem imposes a power-law temporal surface temperature boundary condition at the slab edge ($x = 0$):
\begin{equation}
    T_s(t) = T_0 t^s = 1.0 \times \left(\frac{t}{\text{ns}}\right)^{10/39} \,\text{HeV}
\end{equation}
where the temporal drive exponent is precisely the Henyey exponent $s = s_H$:
\begin{equation}
    s_H = \frac{1}{4 + \alpha - b} = \frac{1}{4 + 1.5 - 1.6} = \frac{10}{39} \approx 0.256410
\end{equation}

For implementations running on an incident radiation field boundary condition (such as Implicit Monte Carlo or transport codes utilizing a Marshak/Milne boundary interface), the corresponding source bath temperature must be explicitly driven over time via:
\begin{equation}
    T_{\text{bath}}(t) = T_s(t) \left[ 1 + \mathcal{B} \right]^{1/4} = \left(\frac{t}{\text{ns}}\right)^{10/39} \left[ 1 + 0.0576603 \left(\frac{t}{\text{ns}}\right)^{-8/13} \right]^{1/4} \,\text{HeV}
\end{equation}
This formulation yields an exact dimensionless boundary flux normalization value $S(0) = 2.003$ matching the internal solver metrics, which tracks the net energy flux flowing through the origin boundary.

\subsubsection{Analytical Wavefront Verification Metrics}
Because $s = s_H$, the similarity profile propagates at a perfectly constant wavefront velocity ($d = 1$), meaning the heat front position $x_F(t)$ changes linearly in time. Code developers should verify their spatial tracking arrays against the following baseline metrics at the final simulation time $t_{\text{final}} = 1\,\text{ns}$:
\begin{itemize}
    \item \textbf{Heat Front Position:} 
    \begin{equation}
        xF(t) = 0.00118584 \times \left(\frac{t}{\text{ns}}\right) \,\text{cm}
    \end{equation}
    \item \textbf{Total System Energy Integral:} 
    \begin{equation}
        E_{\text{total}}(t) = 8.40870 \times 10^9 \times \left(\frac{t}{\text{ns}}\right)^{55/39} \,\text{erg}
    \end{equation}
    \item \textbf{System Optical Thickness:} At $t = 1\,\text{ns}$, the cumulative optical depth across the heated domain scales to $\tau \approx 11.86$, proving the medium is sufficiently thick to validate the LTE diffusion equations against higher-order structural transport solvers.
\end{itemize}

\subsection{Test 2 - Linear Front Wave (High-Temperature Regime)}
\textit{[Content for Test 2 follows here...]}

\end{document}